\chapter*{Abstract}
\addcontentsline{toc}{chapter}{Abstract}

Health information written for a general audience is often too dense for the people who need it most. Patient leaflets, treatment summaries, and public health updates frequently carry the phrasing of a clinical abstract rather than the plain language a worried reader can act on. This project asks whether a transformer-based language model can close some of that gap: can it rewrite technical biomedical text into something easier to read, without quietly dropping the details that make the information medically useful?

The work is built around Cochrane-auto, a dataset that pairs sentences from Cochrane systematic review abstracts with their official plain-language summaries. A small transformer, T5-small, was fine-tuned on these pairs and compared against two baselines: a hand-written rule-based simplifier and the same T5-small model used without any fine-tuning. All three systems were evaluated on a held-out test sample using established readability formulas (Flesch Reading Ease, Flesch-Kincaid Grade Level, SMOG), text generation metrics (BLEU, ROUGE), a simplification-specific metric (SARI), and a semantic similarity measure (BERTScore), followed by a manual reading of individual outputs.

Fine-tuning clearly helped relative to the untrained model. The zero-shot T5-small often produced disfluent or even non-English text, while the fine-tuned version reliably returned coherent sentences and scored well ahead of it on every metric. Against the simpler rule-based baseline the picture was more mixed: the fine-tuned model edged ahead on SARI, ROUGE, and BERTScore, but the rule-based system, and even the untouched source text, sometimes scored as more readable by the surface-level formulas. Manual inspection suggested why. At the training scale used here, the model tends to make small, safe edits, such as removing a statistical parenthetical, rather than the kind of restructuring a human writer produces when a sentence needs deeper rephrasing. In several cases the reference summary answers a slightly different question than the source sentence poses, which is itself a property of the dataset worth flagging for anyone building on this work.

The result is a working, reproducible simplification pipeline rather than a finished product, and the thesis is explicit about that. It documents where automation already helps, where it currently falls short, and why a health-communication tool built this way still needs a human in the loop before it goes near a real patient.
